\begin{proof}
\begin{enumerate}[label=\arabic*.]
\item We remove the existential free-lift condition by an explicit
construction. Let $d>1$ be the order of $M$ and choose a primitive
$u\in\ker(M-I)\cap\mathbb Z^6$. Such a vector is obtained from a
nonzero rational kernel vector by clearing denominators and dividing
by the greatest common divisor of its integer coordinates; both
operations retain the equation $(M-I)u=0$. Put
$g(x)=Mx+u/d^2$ and $S_d=\sum_{a=0}^{d-1}M^a$.
Then $g^j(x)=M^jx+ju/d^2$, $S_dM^j=S_d$ and $S_du=du$.
A fixed point of $g^j$ gives
$(M^j-I)x+ju/d^2=\lambda\in\mathbb Z^6$.
Applying $S_d$ gives $ju/d=S_d\lambda\in\mathbb Z^6$.
A Bezout row pairing to one with $u$ therefore gives $d\mid j$.
Writing $j=dr$, the original equation becomes $ru/d=\lambda$,
so the same row gives $d\mid r$. Thus only multiples of $d^2$
can have fixed points. Since $g^{d^2}$ is translation by the
integer vector $u$, its torus action is the identity. This proves
both its actual order $d^2$ and freeness of every nonidentity power.
For $M=I$, translation by $u/2$ for primitive $u$ has actual order
two and is free. Conversely, when $I-M$ is invertible over
$\mathbb R$, every affine map $Mx+a$ fixes the torus class of
$(I-M)^{-1}a$. The free-lift test is consequently exactly
$\ker(M-I)\ne0$.

\item Finite order is also effectively bounded at this linear
stage. Its minimal polynomial is a product of distinct
cyclotomic polynomials, and eigenvalue one uses at least one
of the six dimensions. Every other cyclotomic factor therefore
has degree at most five and their total degree is at most
five. For a cyclotomic index $d$, each prime divisor $p$ has
$p-1\le5$, and its prime-power factors satisfy
$p^{a-1}(p-1)\le5$. Enumerating these factors and the possible
degree sums gives exactly
\[
                 d\in\{1,2,3,4,5,6,8,10,12\}.             \tag{F.CN1}\label{maps:eq:F.CN1}
\]
Here $d$ is the order of the matrix, the least common multiple
of its cyclotomic indices. Multiplicities of a factor can only
increase the dimension and introduce no new orders. Each
listed order occurs with a fixed vector after adding trivial
summands. Thus taking the powers in \eqref{maps:eq:F.CN1}, checking the
fixed kernel, the product relation, and the cusp Smith
factors gives an exact admissibility test. The unbounded
actual affine multiplicities are not bounded by \eqref{maps:eq:F.CN1}.

\item For an admissible full tuple $T=(M_1,\ldots,M_k,M_0)$ let
$H(T)$ be its maximum absolute entry. There is at most one identity
finite slot in a three-point tuple: two such identities would make
$M_0$ finite order, whereas $M_0=I+N$, $N^2=0$ and
$\operatorname{rank}N=3$ give $(I+N)^d=I+dN\ne I$ for every
positive integer $d$. At a fixed identity slot, delete that matrix
and reinsert $I$ at the same slot after classifying the remaining
ordered tuple. These two maps are inverse, and both commute with
$T\mapsto (QM_1Q^{-1},\ldots,QM_0Q^{-1})$ for the same
$Q\in\mathrm{GL}(6,\mathbb Z)$. They therefore retain the actual
returned simultaneous conjugator. The remaining invertible integral
matrices have height at least one, so inserting $I$ leaves height
unchanged. Every entry in the inserted block is constant; the first
unequal row-major entry of two tuples is therefore the same before
and after insertion. This proves preservation of the specified
lexicographic comparison as well. The actual affine order and
numerator at that slot are retained as gluing data and are transported
by this same $Q$. In the two-involution sector use the canonical
representative of \ref{maps:src:F.2.38}, with its own specified
convention. In all remaining sectors, order full integer tuples
first by $H$, then lexicographically by their row-major entries in
the displayed marked order, and take the first admissible member
of the simultaneous conjugacy class.

\item To compute it, enumerate heights $1,\ldots,H(T)$ and the
integer entries at that height. There are at most
\[
                     (2H(T)+1)^{36(k+1)}                  \tag{F.CN2}\label{maps:eq:F.CN2}
\]
tuples in this complete box. The product relation allows
$M_0$ to be calculated from the first $k$ matrices, so the
implementation enumerates fewer entries. It also rejects
matrices with different characteristic polynomials or orders
before calling the conjugacy decision \ref{maps:src:F.2.36}. These filters
cannot reject an equivalent tuple. At least one comparison
succeeds: the supplied tuple itself belongs to its own box.
The first successful comparison returns the canonical
representative and its exact conjugating matrix.

\item Every simultaneous orbit has a least element in this order.
Its height is at most the height of any supplied member,
which proves that \eqref{maps:eq:F.CN2} loses no class representative for
that input. If two inputs are equivalent they therefore
return the same least element. Conversely the returned
conjugators show that equal representatives imply equivalent
inputs. The casewise involution normal form has the same
property by \ref{maps:src:F.2.38}, and an involution tuple cannot be
equivalent to a tuple in the other branch. This proves the
classification assertion, including its finite stopping rule.

\item Finally, \ref{maps:src:F.2.16} already gives an unbounded admissible integral
family. In its homology convention direct symbolic
multiplication gives
\[
             \operatorname{tr}(M_1M_2^3)
                       =-v^4+v^3-v^2+2.                 \tag{F.CN3}\label{maps:eq:F.CN3}
\]
This is unbounded even on the previously constructed free
affine subfamily $v\in20\mathbb Z$. It is invariant under
simultaneous conjugacy. A fixed finite entry box contains
only finitely many tuples and hence only finitely many such
traces; it cannot contain a representative of every class.
The input-dependent nature of \eqref{maps:eq:F.CN2} is therefore necessary
for this particular exhaustive procedure.

\end{enumerate}
\end{proof}
